% Modelo de trabalho acadêmico (Teses, Dissertações, TCC)
% Documento principal
%
% Universidade Federal do Maranhão - UFMA
% Autor: Sidney Cerqueira <cerqueirasidney@gmail.com>
% Modelo modificado do Cefet-MG
%
% Informações:
%   Codificação utilizada: UTF-8
%   Tamanho da tabulação: 4 (espaços)
\documentclass[oneside]{abntex2-cefetmg}            % Imprimir apenas frente
%\documentclass[doubleside]{abntex2-cefetmg}        % Imprimir frente e verso

% Importações de pacotes
%\usepackage[alf, abnt-emphasize=bf, bibjustif, recuo=0cm, abnt-etal-cite=2, abnt-etal-list=0]{abntex2cite}  % Citações padrão ABNT
\bibliographystyle{ieeetr}
%\usepackage{lmodern}
\usepackage[utf8]{inputenc}                         % Acentuação direta
\usepackage[T1]{fontenc}                            % Codificação da fonte em 8 bits
\usepackage{graphicx}                               % Inserir figuras
\usepackage{amsfonts, amssymb, amsmath}             % Fonte e símbolos matemáticos
\usepackage{booktabs}                               % Comandos para tabelas
\usepackage{verbatim}                               % Texto é interpretado como escrito no documento
\usepackage{multirow, array}                        % Múltiplas linhas e colunas em tabelas
\usepackage{indentfirst}                            % Indenta o primeiro parágrafo de cada seção.
\usepackage{microtype}                              % Para melhorias de justificação?
\usepackage{palatino}                               % Usa a fonte Palatino
\usepackage[algoruled, portuguese]{algorithm2e}     % Escrever algoritmos
\usepackage{float}                                  % Utilizado para criação de floats
%\usepackage[bottom]{footmisc}                      % Mantém as notas de rodapé sempre na mesma posição
%\usepackage{times}                                 % Usa a fonte Times
%\usepackage{lmodern}                               % Usa a fonte Latin Modern
\usepackage{subfig}                                % Posicionamento de figuras
%\usepackage{scalefnt}                              % Permite redimensionar tamanho da fonte
%usepackage{color, colortbl}                       % Comandos de cores
%\usepackage{lscape}                                % Permite páginas em modo "paisagem"
%\usepackage{ae, aecompl}                           % Fontes de alta qualidade
%\usepackage{picinpar}                              % Dispor imagens em parágrafos
\usepackage{latexsym}                              % Símbolos matemáticos
%\usepackage{upgreek}                               % Fonte letras gregas
\usepackage[table,xcdraw]{xcolor}
%\usepackage{subfigure}                              % Pacote para colocar figuras lado a lado
% Definição de comando para colocar dados alinhados a esquerda em tabelas largas com tamanhos variados
\newcolumntype{C}[1]{>{\let\newline\\\arraybackslash\hspace{0pt}}m{#1}}
% Definição de comando para centralizar dados em tabelas largas
\newcolumntype{K}[1]{>{\centering\let\newline\\\arraybackslash\hspace{0pt}}m{#1}}
\newcommand{\tabitem}{~~\llap{\textbullet}~~}    % comando para colocar bullet em item dentro de uma tabela

\newcolumntype{L}{>{\centering\arraybackslash}m{3cm}} % Definindo o tamanho de coluna em tabela

% Inclui o preâmbulo do documento
%
% Documento: Preâmbulo
%

\titulo{Avaliação do Banco de Dados NoSQL HBase no Contexto de Dados Georreferenciados}
%\title{Title in English}
%\subtitulo{Subtítulo do trabalho}
\autor{Rodrigo Santos Balieiro}
\local{Macapá}
\data{Fevereiro de 2019}
\instituicao{Universidade Federal do Amapá
\\ Bacharelado em Ciência da Computação}
\programa{Departamento de Ciências Exatas e Tecnológicas}
\tipotrabalho{Trabalho de Conclusão de Curso}
\preambulo{Trabalho de Conclusão de Curso apresentado à Universidade Federal do Amapá como requisito parcial para obtenção do título de Bacharel em Ciência da Computação.}
\orientador{Thiago Pinheiro do Nascimento}
%\orientador[Orientadora:]{Nome da orientadora}
\coorientador{Dr. Denivaldo Lopes Pavão}
%\coorientador[Coorientadora:]{Nome da coorientadora}
%\linhapesquisa{Sistemas Inteligentes}


% Define as cores dos links e informações do PDF
\makeatletter
\hypersetup{
    portuguese,
    colorlinks,
    linkcolor=black,
    citecolor=black,
    filecolor=blue,
    urlcolor=black,
    breaklinks=true,
    pdftitle={\@title},
    pdfauthor={\@author},
    pdfsubject={\imprimirpreambulo},
    pdfkeywords={abnt, latex, abntex, abntex2}
}
\makeatother

% Redefinição de labels
\renewcommand{\algorithmautorefname}{Algoritmo}
\def\equationautorefname~#1\null{Equa\c c\~ao~(#1)\null}

% Cria o índice remissivo
\makeindex

% Início do documento
\begin{document}

    % Retira espaço extra obsoleto entre as frases.
    \frenchspacing

    % Elementos pré textuais
    \pretextual
    %
% Documento: Capa
%

\makeatletter
\begin{capa}

    \hspace{-2.0cm}
    \begin{minipage}{0.19\textwidth}
        \includegraphics[width=0.8\textwidth]{./04-figuras/unifap.jpg}
    \end{minipage}
    \quad
    \hspace{-1.5cm}
    \begin{minipage}{.9\textwidth}
        \begin{center}
        \normalfont\scshape{\imprimirinstituicao}\\
        \normalfont\scshape{\imprimirprograma}\\
        \abntex@ifnotempty{\imprimirareaconcentracao}
        {%
            \normalfont\scshape{\imprimirareaconcentracao}
        }
        \end{center}
    \end{minipage}

    \vspace*{200pt}

    \begin{center}
        \ABNTEXchapterfont\Large\scshape\imprimirtitulo
        \abntex@ifnotempty{\imprimirsubtitulo}{%
            {\ABNTEXchapterfont\Large\scshape: }{\ABNTEXchapterfont\large\scshape\imprimirsubtitulo}
        }
    \end{center}

    \vspace*{80pt}

    \begin{center}
        \large\normalfont\scshape\textbf\imprimirautor
    \end{center}

    \vspace*{10pt}

    \begin{center}
        \small\imprimirorientadorRotulo{} \imprimirorientador \\
        \abntex@ifnotempty{\imprimircoorientador}
        {%
            \begin{SingleSpacing}\par\end{SingleSpacing}
            
        }
    \end{center}

    \vspace*{\fill}

    \begin{center}
        \normalfont\scshape{\imprimirlocal}\\
        \normalfont\scshape{\imprimirdata}
    \end{center}

\end{capa}
\makeatother
              % Capa
    %
% Documento: Folha de rosto
%

\makeatletter
\begin{folhaderosto}

    \begin{center}
        {\large\normalfont\scshape\textbf\imprimirautor}
    \end{center}

    \vspace*{150pt}

    \begin{center}
        \ABNTEXchapterfont\Large\scshape\imprimirtitulo
        \abntex@ifnotempty{\imprimirsubtitulo}{%
            {\ABNTEXchapterfont\Large\scshape: }{\ABNTEXchapterfont\large\scshape\imprimirsubtitulo}
        }
    \end{center}

    \vspace*{90pt}

    \abntex@ifnotempty{\imprimirpreambulo}{%
        \SingleSpacing
        \begin{tabular}{p{.24\textwidth}p{.15\textwidth}p{.44\textwidth}}
            & \multicolumn{2}{p{.65\textwidth}}{\small\hyphenpenalty=10000{\imprimirpreambulo}} \\ & & \\
            \abntex@ifnotempty{\imprimirareaconcentracao}
            {%
                & \multicolumn{2}{p{.6\textwidth}}{\small\hyphenpenalty=10000{\imprimirareaconcentracaoRotulo\imprimirareaconcentracao}} \\ & & \\
            }
            \abntex@ifnotempty{\imprimirlinhapesquisa}
            {%
                & \multicolumn{2}{p{.6\textwidth}}{\small\hyphenpenalty=10000{\imprimirlinhapesquisaRotulo\imprimirlinhapesquisa}} \\ & & \\
            }
            & \small\imprimirorientadorRotulo & \imprimirorientador
             \\
        \end{tabular}
    }

    \vspace*{\fill}

%    \begin{center}
 %       \normalfont\scshape{\imprimirinstituicao}\\
  %      \normalfont\scshape{\imprimirprograma}\\
   %     \normalfont\scshape{\imprimirlocal}\\
    %    \normalfont\scshape{\imprimirdata}
    %\end{center}
 \begin{center}
        \normalfont\scshape{\imprimirlocal}\\
        \normalfont\scshape{\imprimirdata}
    \end{center}

\end{folhaderosto}
\makeatother
        % Folha de rosto
    %
% Documento: Folha de aprovação
%

\makeatletter
\begin{folhadeaprovacao}

    \begin{center}
        {\large\normalfont\scshape\textbf\imprimirautor}
    \end{center}

    \vspace*{20pt}

    \begin{center}
        \ABNTEXchapterfont\Large\scshape\imprimirtitulo
        \abntex@ifnotempty{\imprimirsubtitulo}{%
            {\ABNTEXchapterfont\Large\scshape: }{\ABNTEXchapterfont\large\scshape\imprimirsubtitulo}
        }
    \end{center}

    \vspace*{15pt}

    \abntex@ifnotempty{\imprimirpreambulo}{%
        \SingleSpacing
        \begin{tabular}{p{.24\textwidth}p{.15\textwidth}p{.44\textwidth}}
            & \multicolumn{2}{p{.65\textwidth}}{\small\hyphenpenalty=10000{\imprimirpreambulo}} \\ & & \\
        \end{tabular}
    }

    \vspace*{70pt}

    \begin{center}
        Trabalho aprovado. \imprimirlocal, $\qquad$ de $\qquad \qquad$ de 2018.
    \end{center}
	
	\vspace*{20pt}
    \begin{center}
        \assinatura{\textbf{\imprimirorientador} \\ Orientador}
        \assinatura{\textbf{Dr. José Walter Cárdenas Sotil} \\ Universidade Federal do Amapá}
        \assinatura{\textbf{Me. Marco Antônio Leal da Silva} \\ Universidade Federal do Amapá}
        %\assinatura{\textbf{Professor} \\ Convidado 3}
        %\assinatura{\textbf{Professor} \\ Convidado 4}
    \end{center}

    \vspace*{\fill}

%    \begin{center}
 %       \normalfont\scshape{\imprimirinstituicao}\\
  %      \normalfont\scshape{\imprimirprograma}\\
   %     \normalfont\scshape{\imprimirlocal}\\
   %     \normalfont\scshape{\imprimirdata}
   % \end{center}
 \begin{center}
        \normalfont\scshape{\imprimirlocal}\\
        \normalfont\scshape{\imprimirdata}
    \end{center}

\end{folhadeaprovacao}
\makeatother
    % Folha de aprovação
	%
% Documento: Dedicatória
%

\begin{dedicatoria}

Dedico este trabalho a comunidade acadêmica de Ciência da Computação, a qual deixo esta singela contribuição bibliográfica.

Em especial a minha esposa Aline e meu filho Paulo Henrique. 

\end{dedicatoria}
       % Dedicatória
	%
% Documento: Agradecimentos
%

\begin{agradecimentos}

Inicialmente, agradeço à Deus, pela vida, saúde e livramentos que até aqui me mantiveram para o desenvolvimento deste trabalho...

A Família pelo apoio, carinho e sacrifício nas diversas noites em que me ausentei para as pesquisas...

Ao Orientador por abraçar a causa e dedicar-se junto a mim nos mais diversos desafios que foram proporcionados pelo tema...

Aos Amigos que sempre estiveram presentes nas adversidades e alegrias...

À Universidade Federal do Amapá -- UNIFAP, ...

Finalmente, agradeço à todos que contribuíram de forma direta ou indireta para a construção desse trabalho. 

\end{agradecimentos}
    % Agradecimentos
    %
% Documento: Epígrafe
%

\begin{epigrafe}

\textit{“O perigo de verdade não é que computadores passem a pensar como humanos, mas sim que humanos passem a pensar como computadores”. -- Sydney Harris}

\end{epigrafe}
          % Epígrafe
    %
% Documento: Resumo (Português)
%

\begin{resumo}
	
Aplicações de georreferenciamento em tempo real usualmente armazenam e gerenciam grandes quantidades de dados produzidos em baixa latência de tempo. Essa situação é desafiadora quando essas aplicações são orientadas a banco de dados relacionais, que possuem baixa escalabilidade para armazenarem e gerenciarem excessivos volumes de dados voláteis e não-voláteis. Este trabalho propõe uma alternativa para esse contexto. 	
	
\textbf{Palavras-chave}: Georreferenciamento. Escalabilidade. Banco de Dados.

\end{resumo}
          % Resumo na língua vernácula
    %
% Documento: Resumo (Inglês)
%

\begin{resumo}[Abstract]

Application of real time georeferencing are traditionally stored and managed in large amountsofdatainlowlatencytime.Thissessioniseffectivewhentheseapplicationsare oriented to relational databases, which have low scalability to store and manage large volumes of volatile and non volatile data. This work is an alternative to this context.

\textbf{Keywords}: Georeferencing. Scalability. Database.

\end{resumo}
          % Resumo em língua estrangeira
    %
% Documento: Lista de figuras
%

\pdfbookmark[0]{\listfigurename}{lof}
\listoffigures*
\cleardoublepage
      % Lista de figuras
    %%
% Documento: Lista de tabelas
%

\pdfbookmark[0]{\listtablename}{lot}
\listoftables*
\cleardoublepage
      % Lista de tabelas
    %%
% Documento: Lista de quadros
%

\pdfbookmark[0]{\listofquadrosname}{loq}
\listofquadros*
\cleardoublepage
      % Lista de quadros
    %%
% Documento: Lista de algoritmos
%

\newcommand{\algoritmoname}{Algoritmo}
\renewcommand{\listalgorithmcfname}{Lista de Algoritmos}

\floatname{algocf}{\algoritmoname}
\newlistof{listofalgoritmos}{loa}{\listalgoritmoname}
\newlistentry{algocf}{loa}{0}

\counterwithout{algocf}{chapter}
\renewcommand{\cftalgocfname}{\algoritmoname\space}
\renewcommand*{\cftalgocfaftersnum}{\hfill--\hfill}

\pdfbookmark[0]{\listalgorithmcfname}{loa}
\listofalgorithms
\cleardoublepage
   % Lista de algoritmos
    %
% Documento: Lista de abreviaturas e siglas
%

\begin{siglas}
    \item[IBM] International Business Machines - Empresa Multinacional de Informática.
    \item[IOT] Internet Of Things - Internet das Coisas.
    \item[MQTT] Message Queuing Telemetry Transport - Protocolo de Telemetria e transporte de mensagens.
    \item[API] Application Programming Interface - Interface de Programação de Aplicativos.
    \item[Javascript] Linguagem de Programação fracamente tipada voltada para Sistemas Web.
    \item[WEB] Sigla em inglês para referir-se a rede mundial de computadores.
    \item[SQL] Service Query Langage - Linguagem Específica de Programação para Banco de Dados
    
\end{siglas}
       % thiago - Lista de abreviaturas e siglas
    %%
% Documento: Lista de símbolos
%

\begin{simbolos}
    \item[$ \Gamma $] Letra grega Gama
    \item[$ \lambda $] Comprimento de ondada
    \item[$ \in $] Pertence
\end{simbolos}
     % Lista de símbolos
    %
% Documento: Sumário
%

\pdfbookmark[0]{\contentsname}{toc}
\tableofcontents*
\cleardoublepage           % Sumário

    % Elementos textuais
    \textual
   	%
% Documento: Introdução
%

\chapter{Introdução}\label{chap:introducao}  

O presente trabalho aborda o contexto dos sistemas de georreferenciamento de tempo real, que, usualmente necessitam de abordagens escaláveis para o gerenciamento e o armazenamento de grandes quantidades de dados, mas que tradicionalmente são implementados para funcionarem sobre sistemas de bancos de dados relacionais, que dispõem de pouca escalabilidade para armazenar e gerenciar excessivos volumes de dados produzidos em baixa latência de tempo. Para contornar essa situação, esta pesquisa visa apresentar uma solução viável para essa problemática, onde ela será utilizada para a geolocalização de viatura da Guarda Civil Municipal de Macapá. Nesse trabalho, os aspectos essenciais da pesquisa proposta são apresentados e discutidos, tais como: a problemática a ser solucionada, os objetivos e justificativas, a metodologia e o cronograma de atividades a serem realizadas e desenvolvidas. Ao final desse capítulo introdutório, também é apresentada a organização do trabalho proposto.

\section {Contexto}

Muitos sistemas de georreferenciamento de tempo real utilizam os conceitos de \textit{"Internet of Things"}, de forma a possibilitar que aparelhos eletrônicos compartilhem dados de localização \cite{Vega}. Esse compartilhamento de dados georreferenciados vem sendo utilizado em diversos domínios de aplicações, que vão desde as situações de desastes naturais, onde usuários envolvidos em operações de proteção civíl necessitam de graus eficientes e eficazes de coordenação \cite{Barroso}; até para medição remota, armazenamento, processamento e visualização de informações químicas \cite{Vega}.

\section{Problemática}

Como este trabalho conterá experimentação de aplicações que envolvem Banco de Dados, \textit{Big data} e Georreferenciamento é importante destacar que:

Na sociedade atual são crescentes os dispositivos que se conectam à internet e dispõem informações que podem ser usadas e podem ser aproveitadas para diversas aplicações computacionais\cite{Vega}. Dentre elas existem as que apresentam informações e cálculos de lugares, distâncias, clima, altura, velocidade, em que é muito comum o desenvolvimento de Sistemas de Informações Geográficos (SIG’s)\cite{Vega}. Esses sistemas são capazes de armazenar dados cartográficos, censitários, cadastros urbano e rural e imagens de satélite, dispondo de mecanismos para tratar as informações, assim como para consultar, recuperar, visualizar e plotar o material na base de dados georreferenciados\cite{Vega}.

Geralmente os dados gerados por Sistemas de Informações Geográficas ou Geoespaciais são dinâmicos, pois para se obter precisão é necessário que esses dados sejam atualizados em tempo real, acumulando um grande volume de dados para serem armazenados em um banco de dados do tipo relacional\cite{Prikh}.

Por esse motivo, faz-se necessário experimentar outro tipo de banco de dados, em que seja suportado um grande fluxo, com maior escalabilidade e tempo de resposta na transição de informações georreferenciais.

\section{Objetivos}
Os objetivos do presente trabalho de pesquisa são apresentados em torno de um único objetivo geral e quatro objetivos específicos.

\subsection{Objetivo Geral}

Avaliar o sistema de bancos de dados não relacional orientado a coluna HBase no armazenamento e no gerenciamento de volumes de dados georreferenciados de tempo real, no contexto da geolocalização de viaturas da Guarda Municipal de Macapá.

\subsection {Objetivos Específicos}

O presente trabalho de pesquisa tem os seguintes objetivos específicos:

\begin{itemize}
	\item Estudar o contexto dos sistemas de geolocalização em tempo real.
	\item Avaliar sistemas de banco de dados escaláveis para grandes volumes de dados.
	\item Propor um estudo de caso para avaliar e comparar banco de dados relacionais e banco de dados não-relacional HBase para armazenarem e gerenciarem excessivos volumes de dados georreferenciados de tempo real, produzidos em baixa latência.
	 
\end{itemize}

\section {Hipótese}

HBase é um sistema de banco de dados não relacional orientado a coluna, capaz de fornecer escalabilidade ilimitada e vertical;e mecanismos eficientes de gerenciamento de grandes volumes de dados georreferenciados de tempo real.

\section {Justificativa}
Muitos estudiosos no Brasil têm dado destaque à Segurança Pública como um dos pilares da Ordem Social. Mais especificamente, em Macapá, de acordo com dados do Fórum Brasileiro de Segurança Pública \cite{forumdeseguranca:2018}, a capital amapaense possui um dos piores cenários com relação à Segurança e os órgãos públicos envolvidos necessitam cada vez mais de investimentos para mudar esse cenário. Um dos investimentos que se pode fazer é na tecnologia computacional para se obter estratégias preventivas no combate ao crime das áreas urbanas.

Isso porque, de acordo com o Instituto Brasileiro de Geografia e Estatística, três em cada quatro domicílios do país possuem acesso à internet\cite{IBGE-2018}.

Dentre as diversas ferramentas tecnológicas desenvolvidas para auxiliar no combate da criminalidade, a Internet Of Things -- IOT, surge com uma grande diversidade de opções que podem ser adotadas pelos agentes de Segurança Pública, principalmente no que se refere a registro de informações em tempo real e tomada de decisões com base em sistemas ágeis.

Aliada a internet das coisas, o registro de informações em banco de dados que possam obter um maior desempenho com um grande fluxo de dados em tempo real pode ser desempenhado com a tecnologia da Apache HBase, pois:

O Apache HBase é um banco de dados de código aberto NoSQL que fornece acesso de leitura/gravação em tempo real a esses grandes conjuntos de dados\cite{CETAX}

Dentre metodologia a ser aplicada e as respostas adquiridas na coleta de dados, pode-se apresentar uma possível ferramenta na contribuição dos esforços das forças de Segurança Pública para o combate à criminalidade e resolução de conflitos nas áreas urbanas de Macapá.
 
Pelo motivo supracitado, apresenta-se no presente trabalho uma abordagem computacional, utilizando tecnologias correlacionadas de Internet Das Coisas, Georreferenciamento e \textit{Big Data} a fim de experimentar ferramentas digitais para uso das forças de Segurança Pública, no caso aqui abordado, a Guarda Civil Municipal de Macapá.

\section {Metodologia}
Para o desenvolvimento e experimentação do trabalho proposto serão utilizados procedimentos metodológicos divididos em duas etapas. O referencial bibliográfico e a experimentação científica.

O referencial bibliográfico (citar autor) destaca não somente a pesquisa textual, mas a base para argumentação, além do desenvolvimento dos códigos a serem usados nos aplicativos dos testes computacionais.

Já a experimentação científica (citar autor) tem como proposta os testes e a coleta de dados para que depois sejam levantadas conclusões com base em análises objetivas ou subjetivas feitas por meio de números ou percentuais.

Como base para o trabalho foi utilizado o artigo, da empresa IBM\cite{IBM}, que é um tutorial para o desenvolvimento de um aplicativo para celular. Tal aplicação usa o GPS do aparelho para enviar informações de geolocalização (latitudes e longitudes) para um ambiente de nuvem da própria IBM.

Esses dados podem ser consumidos por outra ferramenta digital, neste caso, seguindo o artigo de tutorial, os dados foram aproveitados por um sistema \textit{web} escrito em linguagem PHP, usando o protocolo MQTT, para intercomunicação entre as aplicações do hardware do celular, da nuvem e do computador de visualização. Esse esquema pode ser melhor visualizado na Lista de Figuras.

Maiores detalhes desta abordagem serão apresentados nos próximos capítulos.

Vale salientar que, embora haja neste trabalho uma contribuição de escrita no código e modificações nos aplicativos seguidos no tutorial da IBM, o objetivo central aqui proposto não está focado no desenvolvimento da ferramenta, mas na utilidade de sua aplicabilidade. Sendo assim, os resultados a serem alcançados tentarão satisfazer as hipóteses levantadas. 

   	%\chapter{Referencial Teórico}

O presente capítulo apresenta a base teórica fundamental para a estruturação e o desenvolvimento do presente trabalho de pesquisa. Nesse contexto, essa base teórica gira em torno de conceitos relacionados aos sistemas de georreferenciamento de tempo real, \textit{internet of thing}, bancos de dados relacionais e bancos de dados não relacionais.

\section{Sistemas de georreferenciamento de tempo real}

Os sistemas de georreferenciamento ou sistemas de mapas virtuais podem conter duas bases a serem utilizadas: as dinâmicas e as estáticas\cite{MacEachren}. Nas dinâmicas, os pontos de visualização variam de acordo com novas inserções de dados que podem ser cartográficos, envolvendo latitudes e longitudes; Já nas bases de georreferenciamento estáticas, os dados geralmente são representados por imagens ou gráficos armazenados cuja variação representa situações distintas\cite{MacEachren}, como se fossem realizadas várias fotografias em diferentes locais.

Existem, também, outros tipos de dados coletados por sistemas de georreferenciamento, como é o caso....
\section {\textit{Internet of things}}
De acordo com pesquisadores, pode-se definir Internet Das Coisas como um ambiente de objetos físicos interconectados com a internet através de sensores pequenos e embutidos, criando um ecossistema de computação onipresente (ubíqua), voltado para a facilitação do cotidiano das pessoas\cite{Magrani-2018}.

Embora se tenha um entendimento sobre o que é Internet das Coisas, do inglês \textit{Internet Of Things}, com abreviação IoT, existem divergências no que se refere ao conceito. Isso porque para alguns pesquisadores o termo é associado às redes de computadores e aos seus protocolos de comunicação. Para outros a presença da inteligência artificial é fundamental nos dispositivos conectados e, ainda, os que defendem a inter-relação entre as máquinas e pessoas como fundamento do conjunto IoT\cite{Magrani-2018}.

 
\section{Bancos de dados relacionais}
Para entendermos como serão utilizados os registros das informações neste trabalho, será exposto, neste subitem noções de como as informações computacionais são armazenadas e processadas. Para isto devemos apresentar conceitos sobre Dados, Banco de Dados e Sistema de Gerenciamento de Banco de Dados.

Os termos "dado" e "informação" são semelhantes. Certos autores os diferenciam, se referindo a Dado como valores fisicamente registrados no banco de dados, e Informação como sendo o "significado" desses valores para algum usuário \cite{C.Date}.

O significado de "Dado" também pode ser associado a registro, pois é assim que os sistemas de gerenciamento de Banco de Dados tratam as informações de entrada\cite{Kotaro-2005}. Mais adiante iremos abordar com mais detalhes os Sistemas de Gerenciamento de Banco de Dados.

Os Dados precisam ser organizados em um conjunto de informações para serem armazenados e, posteriormente, podem servir para diversos fins como consultas, estatísticas, cálculos, dentre outros. O local em que eles ficam armazenados é conhecido como Banco de Dados\cite{Kotaro-2005}.
\section {NoSQL}

\section{Mineração de Dados e \textit{Big data}}

\chapter{Internet das Coisas - Internet Of Things}

Para entendermos como a Internet Das Coisas foi utilizada nesta pesquisa, primeiramente faremos um apanhado histórico acerca dos principais acontecimentos que resultaram no que conhecemos como as "coisas conectadas" ou dispositivos conectados e, veremos, como essa tecnologia se torna uma tendência para um novo modelo tanto de negócios quanto do cotidiano.

O termo Internet Das Coisas ou "Internet Of Things", em inglês, foi popularmente apresentado, de acordo com relatos por um cientista da Inglaterra chamado Kevin Ashton que, durante uma conferência da Procter \& Grambler, em 1999, falou sobre a troca de dados, em tempo real, por dispositivos, usando a Internet. Mais tarde, em 2009, Ashton escreveu um artigo, intitulado: "A Coisa Internet Das Coisas", em que descreveu como foi sua ideia para descrever o termo que ficou mundialmente conhecido e que é referência para estudos e negócios.

Em seu texto Ashton conta que trabalhava em uma empresa de cosméticos na parte de estoque e ficava chateado de como o setor era impreciso quando se precisavam fazer consultas dos produtos. Ele logo constatou que seria preciso manter um funcionário que ficasse exclusivamente anotando as informações de entrada e saída dos produtos das prateleiras para que tivesse um resultado satisfatório. Porém, isso seria muito custoso para a empresa gerando uma grande perda de tempo e dinheiro só para obter essa precisão de informações. 

Quando Kevin Ashton teve conhecimento sobre as redes sem fio, as RFIDs, teve a ideia de que as "coisas" poderiam se conectar, não mais somente os computadores. E essas coisas ou objetos poderiam transmitir as informações em tempo real, substituindo o trabalho que era de uma pessoa, por um processo automático do objeto, de forma precisa e economizando tempo.


  

\chapter{A Importância das Tecnologias da Computação na Segurança Pública}

Este capítulo abordará questões de como a informática e tecnologias afins têm marcado a dinâmica de trabalho das Áreas de Seguranças desde a inserção de ferramentas digitais até sensores e robôs para a obtenção de melhores resultados.

Para que  possamos discorrer sobre a importância de várias ferramentas computacionais na Segurança Pública, de forma específica, precisamos fazer uma breve contextualização, tanto na atuação das forças de Segurança Pública no Brasil quanto contextualizar também quais novas tecnologias vêm sendo utilizadas como forma de contenção da criminalidade em diversos países através da História. 

\section{Atuação da Segurança Pública no Brasil após o Regime Militar}

Depois que o último presidente militar João Figueiredo, que governou o país, em 1985, o Brasil passou por uma mudança estrutural, sendo promulgada uma nova Constituição, em 1988 e, com ela, novas regras e novos acontecimentos surgiram na sociedade.
A Segurança Pública, de acordo com um artigo publicado por Moema Dutra Freire, em 2012, existem paradigmas na Área Segurança Brasileira que, mesmo com a transição da Ditadura Militar para a Democracia, perduram com iniciativas que chegam a se mostrar ineficientes e ultrapassadas. São três paradigmas:

"Segurança Nacional, vigente durante a Ditadura Militar, Segurança Pública, que se fortalece com a promulgação da Constituição de 1988 e; Segurança Cidadã, perspectiva que tem se ampliado em toda a América Latina e começa a influenciar o debate em segurança no Brasil, a partir de meados de 2000" (FREIRE, 2012).

A autora ainda afirma que esses três paradigmas coexistem, ou seja, o surgimento de um não excluiu o outro e a existência deles perpassa governos, constituições e, até mesmo, ideologias fazendo com que haja uma grande influência no comportamento dos agentes de segurança que, por sua vez, exercem um papel fundamental na "sensação de segurança" divulgada por meios de comunicação, como é o caso de jornais   

\section{Contextualização das ferramentas computacionais na sociedade brasileira}
A presença das redes sociais no cotidiano da população atual é indispensável quando pensamos em fontes de informação, tanto para saber novidades sobre qual time ganhou ontem, quanto para saber o que o presidente dos Estados Unidos afirmou sobre um determinado assunto. Entretanto as redes sociais foram ganhando sua popularidade aos poucos. De acordo com Eric Lassar 1994, as redes sociais apareceram com aplicativos simples de conversação entre grupos fechados ainda por militares norte-americanos, durante o período da Segunda Guerra Mundial, em que o principal objetivo era manter uma comunicação estratégica que não fosse interceptada pelos inimigos. Mais tarde, esses aplicativos chamados de "chat" no inglês, cuja a tradução seria conversa, foram usados por estudantes universitários até que, em 2000, um cientista chamado Mário Porto, escreveu uma obra sobre o fenômeno dos aplicativos de conversação por computadores, o que não somente era troca de mensagens, mas uma rede social.

Com o advento do Facebook, em 2002 e do Twitter, em 2003, as redes sociais ganharam um grande espaço nas rotinas sociais, passando por uma explosão de visualizações que mudaram a forma de como as pessoas buscavam informações e notícias, pois eram muito mais rápidas e selecionadas do que outros meios de comunicação como TV, rádio, jornais impressos e, até mesmo, sites oficiais de notícias, o que desenvolveu problemas de confiabilidade e crimes digitais, que veremos mais adiante. Ainda abordando esse assunto, podemos citar a compra do aplicativo Whatsapp pelo grupo Facebook, em 2006, fato que fez o aplicativo de celular se tornar a principal ferramenta de troca de informações e dados da segunda década do século XXI, de acordo com dados de uma pesquisa apontada por uma revista científica, em 2008, a ComputerScience, que abordou tanto o número de aparelhos celulares, quando os dispositivos SMART conectados e destacou os principais softwares usados nesses dispositivos, dando destaque para o Facebook e o Whatsapp.

Dessa forma podemos inferir que a rede de computadores e o uso de aplicativos de celulares estão cada vez mais presentes no cotidiano e nas áreas que envolvem ciência/pesquisa e negócios, sendo bastante explorada para que possamos analisar com maiores seus efeitos e, mais a frente, abordarmos as principais formas de se utilizar o conjunto entre essas tecnologias como ferramentas do experimento científico que contemplamos neste trabalho.

   


    %
% Documento: Referências Bibliográficas
%
\bibliography{./refbase}    % Geração automática das referências por meio do arquivo 'refbase.bib'

    

\end{document}
